\section{Conclusion}
\label{section-conclusion}

In this paper we have developed the $\minicat$ language for defining
low-level MPC protocols, and the $\metaprot$ metaprogramming language
for specifying $\minicat$ protocols compositionally and at a higher
level of abstraction. We have developed a Hoare logic based on
SMT-style verification for $\minicat$, as well as confidentiality and
integrity type systems that enforce passive and malicious security in
the real/ideal model. These establish a foundation for an automated
Hoare logic and type system in $\metaprot$ that is compositional,
scalable, and generalizes to arbitrary prime fields and secret sharing
schemes.

Our automated $\metaprot$ Hoare logic and type theory leverages
Satisfiability Modulo Finite Fields for efficiency, scalability, and
generalizability to arbitrary prime fields. We developed examples of
real binary and arithmetic circuit protocols, illustrated the
applicability of our analyses to enforce confidentiality and integrity
properties, and evaluated performance of our implementation on these
examples.

\subsection{Future Work}

There are several directions for future work on the topics explored in
this paper through the $\metaprot/\minicat$ framework. One interesting
theoretical question is whether there are conditions under which
verification of a protocol in $\fieldp{2}$ generalizes to arbitrary
prime fields. It is often the case that MPC protocols are presented in
a way that generalizes over the field size, but it is not clear how to
articulate necessary conditions for generalizing from verification in
$\fieldp{2}$. This would have an obvious advantage since verification in
$\fieldp{2}$ is more efficient than verification in larger fields.

With regard to efficiency of implementation, integration of our system
with the Swanky framework \cite{swanky} would be an interesting
direction to explore. This framework could present opportunities for
accelerating SMT verification due to the optimizations especially
in their SAW backend.

Another major direction to explore would be an extension of the
$\metaprot/\minicat$ framework to non-information-theoretic, or
\emph{computational} security. This is security gained through
computaitonal hardness of, e.g., public or symmetric key cryptography.
Several MPC protocols rely on this sort of security, such as versions
of OT, committment, garbled circuits, etc. Making this extension would
require theoretical and practical extensions of $\metaprot/\minicat$
and make it applicable to a broader class of MPC protocols.
